\chapter{Closed points, divisors and skyscraper cohomology}
\label{chap:RiemannRoch}

This chapter assembles the local and divisorial foundations of the Riemann--Roch ledger of a
smooth curve. On an integral scheme \(X\) smooth of relative dimension one over a field \(K\),
every point other than the generic point is closed, and the stalk of the structure sheaf there
is a discrete valuation ring; this yields, on the function field \(K(X)\), the order-of-vanishing
valuation \(\mathrm{ord}_x\) at each closed point \(x\). Each closed point also carries a residue
degree \([\kappa(x) : K]\), the \(K\)-dimension of its residue field. A Weil divisor is then a
finitely supported \(\Z\)-family on the closed points, and its degree is the residue-degree
weighted sum of its coefficients --- the correct notion over a non-closed base field. Finally, to
each \(K\)-module \(M\) and point \(x\) is attached the skyscraper sheaf of \(K\)-modules
concentrated at \(x\); its degree-zero cohomology recovers \(M\) and its degree-one cohomology
vanishes, the keystone input to the local contribution of a point in the cohomological
Riemann--Roch computation.

Throughout, \(K\) is a field, \(X\) is an integral scheme smooth of relative dimension one over
\(\Spec K\) --- the curve bundle of Chapter~\ref{chap:Curves} --- and \(\eta\) denotes its generic
point. We write \(\struct{X,x}\) for the stalk at \(x\), \(\kappa(x)\) for its residue field,
and \(K(X)\) for the function field (the stalk at \(\eta\)).

\section{Discrete valuations at closed points}

\begin{theorem}[Discrete valuation ring at a closed point of an affine Dedekind chart]
  \label{thm:affine_dedekind_dvr_stalk}
  \lean{AlgebraicGeometry.IsAffineOpen.isDiscreteValuationRing_stalk}
  \leanok
  Let \(X\) be an integral scheme, \(V \subseteq X\) an affine open whose ring of sections
  \(\Gamma(X, V)\) is a Dedekind domain, and \(x \in V\) a point which is not the generic point
  of \(X\). Then the stalk \(\struct{X,x}\) is a discrete valuation ring.
\end{theorem}
\begin{proof}
  \leanok
  Under the isomorphism \(V \cong \Spec \Gamma(X, V)\) the point \(x\) corresponds to a prime
  ideal \(\mathfrak q\) of \(\Gamma(X, V)\), and the stalk \(\struct{X,x}\) is the localization
  of \(\Gamma(X, V)\) at \(\mathfrak q\). The prime \(\mathfrak q\) is nonzero: if
  \(\mathfrak q = 0\) then \(x\) is the image of the generic point of \(\Spec \Gamma(X, V)\),
  which is the generic point of the nonempty open \(V\), hence the generic point of the
  irreducible scheme \(X\) --- contrary to hypothesis. The localization of a Dedekind domain at
  a nonzero prime ideal is a discrete valuation ring, which is the claim.
\end{proof}

\begin{theorem}[Discrete valuation ring at a closed point of a curve]
  \label{thm:closedpoint_dvr_stalk}
  \lean{AlgebraicGeometry.isDiscreteValuationRing_stalk}
  \uses{thm:affine_dedekind_dvr_stalk, thm:relDim1_dedekind_charts}
  \leanok
  Let \(X\) be an integral scheme smooth of relative dimension one over a field \(K\), and let
  \(x \in X\) be a point other than the generic point. Then the stalk \(\struct{X,x}\) is a
  discrete valuation ring.
\end{theorem}
\begin{proof}
  \leanok
  By \ref{thm:relDim1_dedekind_charts} the point \(x\) lies in an affine open \(V\) whose ring
  of sections \(\Gamma(X, V)\) is a Dedekind domain. Apply \ref{thm:affine_dedekind_dvr_stalk}.
\end{proof}

\begin{theorem}[Dedekind stalk at a closed point of a curve]
  \label{thm:closedpoint_dedekind_stalk}
  \lean{AlgebraicGeometry.isDedekindDomain_stalk}
  \uses{thm:closedpoint_dvr_stalk}
  \leanok
  In the situation of \ref{thm:closedpoint_dvr_stalk}, the stalk \(\struct{X,x}\) is a Dedekind
  domain.
\end{theorem}
\begin{proof}
  \leanok
  A discrete valuation ring is a local principal ideal domain, in particular a Dedekind domain;
  apply this to the stalk, which is a discrete valuation ring by \ref{thm:closedpoint_dvr_stalk}.
\end{proof}

\begin{theorem}[Non-generic points of a curve are closed]
  \label{thm:closedpoint_isClosed}
  \lean{AlgebraicGeometry.isClosed_singleton_of_ne_genericPoint}
  \uses{thm:relDim1_specializes, lem:height_one_closed_points}
  \leanok
  Let \(X\) be an integral scheme smooth of relative dimension one over a field \(K\), and let
  \(x \in X\) be a point other than the generic point. Then the singleton \(\{x\}\) is closed.
\end{theorem}
\begin{proof}
  \leanok
  By \ref{thm:relDim1_specializes} the specialization order of \(X\) has height one: every
  generalization of any point is the generic point or the point itself. A scheme with
  irreducible underlying space and height-one specialization order has all its non-generic
  points closed (\ref{lem:height_one_closed_points}); apply this to \(x\).
\end{proof}

\begin{definition}[The order valuation at a closed point]
  \label{def:scheme_ord}
  \lean{AlgebraicGeometry.Scheme.ord}
  \uses{thm:closedpoint_dvr_stalk, thm:closedpoint_dedekind_stalk}
  \leanok
  \source{papaioannou-algebraic-rr:page-0003}
  For a closed point \(x\) of \(X\) (a point other than the generic point), the
  \emph{order valuation} \(\mathrm{ord}_x\) is the discrete valuation of the function field
  \(K(X)\) with values in \(\Z \cup \{+\infty\}\) determined by the discrete valuation ring
  \(\struct{X,x}\): it assigns to a nonzero rational function its order of vanishing at \(x\),
  and \(+\infty\) to \(0\).
\end{definition}
\begin{proof}
  \leanok
  The stalk \(\struct{X,x}\) is a discrete valuation ring (\ref{thm:closedpoint_dvr_stalk}),
  hence a Dedekind domain (\ref{thm:closedpoint_dedekind_stalk}) which is not a field, so its
  maximal ideal is a nonzero prime, i.e.\ a height-one prime of \(\struct{X,x}\). On an integral
  scheme the function field \(K(X)\) is the fraction field of every stalk, in particular of
  \(\struct{X,x}\); and a height-one prime of a Dedekind domain \(B\) determines a discrete
  valuation on \(\Frac B\) --- the \(\mathfrak p\)-adic valuation counting the multiplicity of
  \(\mathfrak p\) in the factorization of a fractional ideal. Taking \(B = \struct{X,x}\) and
  \(\mathfrak p\) its maximal ideal gives the order valuation on \(K(X)\).
\end{proof}

\section{The residue degree}

The curve \(X\) is a scheme over \(\Spec K\); as in Chapter~\ref{chap:Cohomology} the base field
acts on the sections through the structure morphism, and the same discipline supplies the module
structure on each residue field.

\begin{definition}[The base-field algebra map on a residue field]
  \label{def:residueOverAlgebraMap}
  \lean{AlgebraicGeometry.Scheme.residueOverAlgebraMap}
  \uses{def:overAlgebraMap}
  \leanok
  For a scheme \(X\) over \(\Spec K\) and a point \(x \in X\), the canonical ring homomorphism
  \[
    K \longrightarrow \kappa(x)
  \]
  through the structure morphism: the composite of the global algebra map
  \(K \to \Gamma(X, \struct X)\) of \ref{def:overAlgebraMap}, the germ
  \(\Gamma(X, \struct X) \to \struct{X,x}\) at \(x\), and the residue map
  \(\struct{X,x} \to \kappa(x)\). By restriction of scalars along it, the residue field
  \(\kappa(x)\) is a \(K\)-module (indeed a \(K\)-algebra); as with the module structure on the
  structure sheaf, this is used as an explicit map rather than a global instance.
\end{definition}
\begin{proof}
  \leanok
  Each of the three maps is a ring homomorphism, and \(K \to \Gamma(X, \struct X)\) is the base
  algebra map of \ref{def:overAlgebraMap}; their composite is the stated ring homomorphism.
\end{proof}

\begin{definition}[The residue degree]
  \label{def:residueDeg}
  \lean{AlgebraicGeometry.Scheme.residueDeg}
  \uses{def:residueOverAlgebraMap}
  \leanok
  \source{papaioannou-algebraic-rr:page-0004}
  For a scheme \(X\) over \(\Spec K\) and a point \(x \in X\), the \emph{residue degree} at
  \(x\) is
  \[
    \deg_K(x) \;=\; [\kappa(x) : K] \;=\; \dim_K \kappa(x),
  \]
  the \(K\)-dimension of the residue field, taken with the \(K\)-module structure of
  \ref{def:residueOverAlgebraMap}.
\end{definition}

\section{Weil divisors and their degree}

A Weil divisor on the curve records a finite prescription of orders at closed points. Its degree
must weight each point by its residue degree over the base field: this is the correct notion when
\(K\) is not algebraically closed. For example, on \(\PP^1_{\mathbb{R}}\) the rational function
\(t^2 + 1\) has a single closed zero \(x\), whose residue field is \(\mathbb{C}\), and a double
pole at \(\infty\), with residue field \(\mathbb{R}\); the unweighted point count \(1 - 2 = -1\)
is nonzero, whereas the residue-degree weighted degree
\(1 \cdot [\mathbb{C} : \mathbb{R}] - 2 \cdot [\mathbb{R} : \mathbb{R}] = 2 - 2 = 0\)
vanishes, as the degree of a principal divisor must.

\begin{definition}[Weil divisors on a curve]
  \label{def:curveDivisor}
  \lean{AlgebraicGeometry.Scheme.CurveDivisor}
  \uses{thm:closedpoint_isClosed}
  \leanok
  \source{papaioannou-algebraic-rr:page-0006}
  For an integral scheme \(X\) with generic point \(\eta\), a \emph{Weil divisor} is a finitely
  supported \(\Z\)-valued function on the non-generic points, i.e.\ an element of the free abelian
  group
  \[
    \Div(X) \;=\; \bigl(\{ x \in X : x \neq \eta \} \to_{\mathrm{fin}} \Z\bigr)
  \]
  on the set of non-generic points. It carries the pointwise abelian-group structure and the
  pointwise partial order, both independent of any base field. For the curve \(X\) the
  non-generic points are exactly the closed points (\ref{thm:closedpoint_isClosed}), so a Weil
  divisor is a finite formal \(\Z\)-combination of closed points.
\end{definition}

\begin{definition}[The degree of a Weil divisor]
  \label{def:curveDivisor_deg}
  \lean{AlgebraicGeometry.Scheme.CurveDivisor.deg}
  \uses{def:curveDivisor, def:residueDeg}
  \leanok
  \source{papaioannou-algebraic-rr:page-0007}
  For a scheme \(X\) over \(\Spec K\) and a Weil divisor \(D \in \Div(X)\), the \emph{degree} of
  \(D\), weighted by the residue degrees over the base field, is
  \[
    \deg_K D \;=\; \sum_{x} D_x \cdot [\kappa(x) : K] \;\in\; \Z ,
  \]
  the finite sum over the support of \(D\) of the coefficients \(D_x\) times the residue degrees
  \([\kappa(x) : K]\) of \ref{def:residueDeg}.
\end{definition}

\begin{theorem}[The degree is a group homomorphism]
  \label{thm:curveDivisor_deg_hom}
  \lean{AlgebraicGeometry.Scheme.CurveDivisor.deg_zero,
    AlgebraicGeometry.Scheme.CurveDivisor.deg_add,
    AlgebraicGeometry.Scheme.CurveDivisor.deg_neg}
  \uses{def:curveDivisor_deg}
  \leanok
  \source{papaioannou-algebraic-rr:page-0007}
  The degree map \(\deg_K : \Div(X) \to \Z\) is a homomorphism of abelian groups: it satisfies
  \(\deg_K 0 = 0\), \(\deg_K(D + D') = \deg_K D + \deg_K D'\) and \(\deg_K(-D) = -\deg_K D\).
\end{theorem}
\begin{proof}
  \leanok
  The degree is the finitely supported sum \(D \mapsto \sum_x D_x \cdot [\kappa(x) : K]\). The
  zero divisor has empty support, so \(\deg_K 0 = 0\). For additivity, the summand
  \(n \mapsto n \cdot [\kappa(x) : K]\) is, at each point \(x\), additive in the coefficient
  \(n\) and vanishes at \(n = 0\); hence the finitely supported sum of the pointwise sums
  \((D + D')_x = D_x + D'_x\) splits as \(\sum_x D_x[\kappa(x):K] + \sum_x D'_x[\kappa(x):K]\),
  giving \(\deg_K(D + D') = \deg_K D + \deg_K D'\). Finally, applying additivity to \(D\) and
  \(-D\) and using \(\deg_K 0 = 0\) yields \(\deg_K(-D) = -\deg_K D\).
\end{proof}

\begin{lemma}[Degree of a single-point divisor]
  \label{lem:curveDivisor_deg_single}
  \lean{AlgebraicGeometry.Scheme.CurveDivisor.deg_single}
  \uses{def:curveDivisor_deg}
  \leanok
  For a closed point \(x\) and an integer \(n\), the degree of the single-point divisor
  \(n \cdot x\) is
  \[
    \deg_K(n \cdot x) \;=\; n \cdot [\kappa(x) : K] .
  \]
\end{lemma}
\begin{proof}
  \leanok
  The divisor \(n \cdot x\) is supported at the single point \(x\) with coefficient \(n\), so the
  defining sum of \ref{def:curveDivisor_deg} has the single term \(n \cdot [\kappa(x) : K]\)
  (the summand vanishes at coefficient \(0\), so no other point contributes).
\end{proof}

\section{Skyscraper sheaves and the vanishing of \(H^1\)}

Fix a scheme \(X\) over \(\Spec K\) and work on its small Zariski site (the poset of opens with
open coverings), whose terminal object is the whole space \(\top = X\). For a point \(x \in X\)
and a \(K\)-module \(M\) the skyscraper sheaf concentrates \(M\) at \(x\).

\begin{definition}[The skyscraper sheaf of \(K\)-modules]
  \label{def:skyModule}
  \lean{AlgebraicGeometry.skyModule}
  \leanok
  For \(x \in X\) and a \(K\)-module \(M\), the \emph{skyscraper sheaf} \(\mathrm{sky}_x(M)\) of
  \(K\)-modules on the small Zariski site of \(X\) is the sheaf whose sections over an open
  \(U\) are
  \[
    \mathrm{sky}_x(M)(U) \;=\;
      \begin{cases} M & x \in U, \\ 0 & x \notin U, \end{cases}
  \]
  with restriction between two opens containing \(x\) the identity of \(M\), and to an open
  avoiding \(x\) the zero map. Its restriction map between any two opens that both contain \(x\)
  is an isomorphism, and its sections over any open avoiding \(x\) are the zero module.
\end{definition}

\begin{definition}[Degree-zero cohomology of a skyscraper is its value]
  \label{def:skyModuleGammaEquiv}
  \lean{AlgebraicGeometry.skyModuleGammaEquiv}
  \uses{def:skyModule, def:HModule_linearEquiv0}
  \leanok
  For \(x \in X\) and a \(K\)-module \(M\), there is a \(K\)-linear equivalence
  \[
    H^0\bigl(X, \mathrm{sky}_x(M)\bigr) \;\cong\; M .
  \]
\end{definition}
\begin{proof}
  \leanok
  Global sections live over the terminal open \(\top = X\), which contains \(x\), so
  \(\mathrm{sky}_x(M)(\top) = M\) by \ref{def:skyModule}. Compose with the identification of
  degree-zero cohomology with global sections, \(H^0(F) \cong F(\top)\) of
  \ref{def:HModule_linearEquiv0}, applied to \(F = \mathrm{sky}_x(M)\).
\end{proof}

\begin{theorem}[Global lifting along the cokernel of a skyscraper embedding]
  \label{thm:sky_cokernel_app_top_surjective}
  \lean{AlgebraicGeometry.cokernelπ_app_top_surjective}
  \uses{def:skyModule, lem:sections_left_exact, lem:app_injective_of_mono}
  \leanok
  Let \(x \in X\), let \(M\) be a \(K\)-module, and let \(\iota : \mathrm{sky}_x(M) \to G\) be a
  monomorphism of sheaves of \(K\)-modules. Then every global section of \(\coker \iota\) lifts
  to a global section of \(G\): the map \(G(\top) \to (\coker \iota)(\top)\) is surjective.
\end{theorem}
\begin{proof}
  \leanok
  Write \(\pi : G \to \coker \iota\) for the projection and let \(q \in (\coker \iota)(\top)\).
  The proof glues corrected local lifts, the correction being an \emph{elementwise} coboundary
  supplied by flasqueness of the skyscraper --- so no quasi-compactness or finite subcover is
  needed. Restriction maps of a sheaf \(F\) are written \(t|_W\).

  (a) \emph{Local lifts.} As an epimorphism of sheaves, \(\pi\) is locally surjective on
  sections: for each point \(p \in X\) there is an open \(W_p \ni p\) and a section
  \(s_p \in G(W_p)\) with \(\pi(s_p) = q|_{W_p}\). The opens \(W_p\) cover \(X\), since each
  contains its point \(p\).

  (b) \emph{The differences come from the skyscraper.} On \(W_p \cap W_{p'}\) the difference
  \(s_p|_{W_p \cap W_{p'}} - s_{p'}|_{W_p \cap W_{p'}}\) maps to
  \(q|_{W_p \cap W_{p'}} - q|_{W_p \cap W_{p'}} = 0\) under \(\pi\). Since \(\iota\) is a
  monomorphism it is the kernel of its cokernel on sections (\ref{lem:sections_left_exact}), so
  there is a unique \(a_{pp'} \in \mathrm{sky}_x(M)(W_p \cap W_{p'})\) with
  \(\iota(a_{pp'}) = s_p - s_{p'}\) there.

  (c) \emph{Cocycle identity.} On a triple overlap \(W_p \cap W_{p'} \cap W_l\),
  \[
    \iota(a_{pl}) = s_p - s_l = (s_p - s_{p'}) + (s_{p'} - s_l) = \iota(a_{pp'} + a_{p'l}) ,
  \]
  and \(\iota\) is injective on sections (\ref{lem:app_injective_of_mono}), so
  \(a_{pl} = a_{pp'} + a_{p'l}\) there.

  (d) \emph{Pivot on the support point.} Because \(\mathrm{sky}_x(M)\) is a skyscraper sheaf it
  is \emph{flasque}: every restriction map of \(\mathrm{sky}_x(M)\) is surjective. In particular
  the restriction \(\mathrm{sky}_x(M)(W_p) \to \mathrm{sky}_x(M)(W_p \cap W_x)\) is surjective,
  so choose \(b_p \in \mathrm{sky}_x(M)(W_p)\) with \(b_p|_{W_p \cap W_x} = a_{px}\), where the
  index \(x\)-th cover member \(W_x \ni x\) is the neighbourhood of the support point itself.

  (e) \emph{Restriction to \(W_x\) is injective.} For every open \(A\), the restriction
  \(\mathrm{sky}_x(M)(A) \to \mathrm{sky}_x(M)(A \cap W_x)\) is injective: if \(x \in A\) then
  both opens contain \(x\) and the restriction is an isomorphism (\ref{def:skyModule}); if
  \(x \notin A\) then \(\mathrm{sky}_x(M)(A) = 0\) is trivial, so any map out of it is injective.

  (f) \emph{The cocycle is the coboundary of \(b\).} We claim
  \(a_{pp'} = b_p|_{W_p \cap W_{p'}} - b_{p'}|_{W_p \cap W_{p'}}\). By (e) it suffices to check
  this after restriction to \((W_p \cap W_{p'}) \cap W_x\). There the cocycle identity (c) with
  \(l = x\) gives \(a_{pp'} = a_{px} - a_{p'x}\), and \(a_{px} = b_p\), \(a_{p'x} = b_{p'}\)
  after restriction to \(\cap W_x\) by the choice in (d); hence
  \(a_{pp'} = b_p - b_{p'}\) restricted to \((W_p \cap W_{p'}) \cap W_x\), as required.

  (g) \emph{Corrected local sections glue.} Put \(t_p = s_p - \iota(b_p) \in G(W_p)\). On
  \(W_p \cap W_{p'}\),
  \[
    t_p - t_{p'} = (s_p - s_{p'}) - \iota(b_p - b_{p'}) = \iota(a_{pp'}) - \iota(a_{pp'}) = 0 ,
  \]
  using (b) and (f). So the \(t_p\) agree on overlaps and, by the sheaf condition, glue to a
  global section \(s \in G(\top)\) with \(s|_{W_p} = t_p\).

  (h) \emph{The glued section lifts \(q\).} On each \(W_p\), since \(\pi \circ \iota = 0\),
  \[
    \pi(s)|_{W_p} = \pi(t_p) = \pi(s_p) - \pi(\iota(b_p)) = q|_{W_p} .
  \]
  Sections of the sheaf \(\coker \iota\) that agree on the cover \(\{W_p\}\) are equal, so
  \(\pi(s) = q\). Thus \(s\) is the required lift.
\end{proof}

\begin{theorem}[Degree-one cohomology of a skyscraper vanishes]
  \label{thm:skyModule_h1_vanishing}
  \lean{AlgebraicGeometry.skyModule_subsingleton_hModule_one}
  \uses{def:skyModule, def:HModule}
  \leanok
  For every point \(x \in X\) and every \(K\)-module \(M\),
  \[
    H^1\bigl(X, \mathrm{sky}_x(M)\bigr) \;=\; 0 .
  \]
\end{theorem}
\begin{proof}
  \leanok
  \uses{lem:ext_one_vanishing_criterion, thm:sky_cokernel_app_top_surjective, def:constHomEquiv,
    lem:constHomEquiv_naturality}
  The category of sheaves of \(K\)-modules on the small Zariski site of \(X\) is Grothendieck
  abelian, hence has enough injectives: embed \(\mathrm{sky}_x(M)\) into an injective sheaf by a
  monomorphism \(\iota : \mathrm{sky}_x(M) \to G\). By the \(\Ext^1\)-vanishing criterion
  \ref{lem:ext_one_vanishing_criterion} (with \(L = K_J\) the constant sheaf), it suffices to
  lift every morphism \(K_J \to \coker \iota\) along \(\pi : G \to \coker \iota\). By
  \ref{def:constHomEquiv} and its naturality \ref{lem:constHomEquiv_naturality}, morphisms out
  of \(K_J\) are global sections, compatibly with postcomposition; so the required lifting is
  exactly the surjectivity of \(G(\top) \to (\coker \iota)(\top)\), which is
  \ref{thm:sky_cokernel_app_top_surjective}.
\end{proof}
